%! Direct, Inverse \& Joint Variation — Unit 5, Lesson 5.7 — Warm-Up
\documentclass[10pt]{article}
\usepackage{algebra2-article}
\usepackage{algebra2-boxes}
\newcommand{\TallMath}[1]{$\displaystyle #1\rule[-1.4em]{0pt}{3.2em}$}

\begin{document}

\pageheader{Unit 5, Lesson 5.7}{Warm-Up}
\namedateperiod

\vspace{0.08in}

\begin{notesbox}{Getting Ready: Skills We'll Use Today}
\small Three skills, all of them yours already. Today is the last lesson of the unit, and it is where
$y=\dfrac{k}{x}$ takes off the lab coat and goes to work.

\vspace{4pt}
\textbf{1. Two rows of arithmetic.} For each table, fill in the \emph{quotient} row and the
\emph{product} row. Decimals are fine.
\begin{center}
\begin{minipage}{0.47\linewidth}\centering
\textbf{Table P}\\[2pt]
\renewcommand{\arraystretch}{1.35}
\begin{tabular}{|l|c|c|c|c|}
\hline
\rowcolor{forestbg}
$x$ & $2$ & $4$ & $6$ & $10$ \\ \hline
$y$ & $10$ & $20$ & $30$ & $50$ \\ \hline
$y\div x$ & & & & \\ \hline
$x\cdot y$ & & & & \\ \hline
\end{tabular}
\end{minipage}\hfill
\begin{minipage}{0.47\linewidth}\centering
\textbf{Table Q}\\[2pt]
\renewcommand{\arraystretch}{1.35}
\begin{tabular}{|l|c|c|c|c|}
\hline
\rowcolor{forestbg}
$x$ & $2$ & $4$ & $6$ & $10$ \\ \hline
$y$ & $30$ & $15$ & $10$ & $6$ \\ \hline
$y\div x$ & & & & \\ \hline
$x\cdot y$ & & & & \\ \hline
\end{tabular}
\end{minipage}
\end{center}
\vspace{-2pt}
In Table P the \blank{2.4cm} row came out constant. In Table Q the \blank{2.4cm} row did.
Hold on to that --- it is the entire test we build today.

\vspace{5pt}
\textbf{2. The rational parent} (Lesson 5.4). Let $y=\dfrac{6}{x}$.
\begin{enumerate}[label=(\alph*), leftmargin=2.1em, itemsep=3pt]
  \item $x=2 \Rightarrow y=\blank{1.2cm}$; \quad $x=3 \Rightarrow y=\blank{1.2cm}$; \quad
        $x=12 \Rightarrow y=\blank{1.2cm}$
  \item Vertical asymptote: \blank{1.8cm} \qquad Horizontal asymptote: \blank{1.8cm}
  \item As $x$ grows without bound, $y$ gets closer and closer to \blank{1.4cm} but never
        \blank{2.6cm} it.
\end{enumerate}

\vspace{5pt}
\textbf{3. Solve the proportion} (Lesson 5.6 --- one fraction $=$ one fraction, so cross-multiply).
\begin{enumerate}[label=(\alph*), leftmargin=2.1em, itemsep=3pt]
  \item $\dfrac{3}{8}=\dfrac{x}{20}$: \quad $x=\blank{2.0cm}$
  \item $\dfrac{5}{x}=\dfrac{2}{9}$: \quad $x=\blank{2.0cm}$
\end{enumerate}
\end{notesbox}

\end{document}
